\documentclass{article}
\usepackage[margin=2cm]{geometry}
\usepackage{amsmath}
\usepackage{graphicx}
\usepackage{booktabs}
\usepackage[utf8]{inputenc}
\begin{document}
This motivates us to define a superpotential \(W\) as
\[
W=\sqrt{S_{0}^{2}+S_{3}^{2}}
\]
Unfortunately, this superpotential does not allow one to write the BPS equations for both \(\phi^{0}\) and \(\phi^{3}\) as gradient flow equations. The reason for this failure is that the integrability condition required to convert the BPS equation into a gradient flow form is not satisfied; see e.g. Appendix C.2.1 of . We thus follow the strategy of to construct an approximate superpotential. Our model consists of two consistent truncations that admit flat domain walls and an exact superpotential. These are the \(\phi^{3}=0, \phi^{0} \neq 0\) truncation and the \(\phi^{0}=0, \phi^{3} \neq 0\) truncation. The corresponding flow equations are (we set \(\eta=-1\) henceforth)
\[
\phi^{0^{\prime}}=-8 \partial_{\phi^{0}} W\left|_{\phi^{3}=0} \quad \phi^{3^{\prime}}=8 \partial_{\phi^{3}} W\right|_{\phi^{0}=0}
\]
respectively. In either truncation, the BPS equations for the warp factor and dilaton \(\sigma\) can be put in the following form,
\[
f^{\prime}=2 W \quad \sigma^{\prime}=2 \partial_{\sigma} W
\]
An important fact is that, though the gradient flow equations of do not hold exactly in the full model with \(\phi^{0} \neq 0, \phi^{3} \neq 0\), they do hold up to and including \(O\left(z^{5}\right)\). Looking at the form of the UV asymptotics of the scalar fields, one may expand the superpotential of keeping only terms contributing up to this order. This gives
\[
W=\frac{1}{2}+\frac{3}{4} \sigma^{2}+\frac{1}{16}\left(\phi^{0}\right)^{2}-\frac{3}{16}\left(\phi^{3}\right)^{2}+\frac{1}{192}\left(\phi^{0}\right)^{4}-\frac{3}{16}\left(\phi^{0}\right)^{2} \sigma+\ldots
\]
where the dots represent terms of order \(O\left(z^{6}\right)\). This is the approximate superpotential we will use in what follows.
\subsection*{0.0.1 Bogomolnyi trick}
We now use the Bogomolnyi trick to get the finite counterterms needed to preserve supersymmetry in the case of a flat domain wall. The central idea of the Bogomolnyi trick is that for a BPS solution, the renormalized on-shell action must vanish. In order to make use of this fact, we will first want to recast the on-shell action in a simpler form. To do so, we begin by inserting. We find that
\[
\mathcal{L}=-\frac{1}{4} R-20 W^{2}+2 \mathcal{L}_{\text {kin }}
\]
where we've defined
\[
\mathcal{L}_{\text {kin }}=\left(\sigma^{\prime}\right)^{2}+\frac{1}{4}\left[-\left(\phi^{3^{\prime}}\right)^{2}+\cos ^{2} \phi^{3}\left(\phi^{0^{\prime}}\right)^{2}\right]
\]
\end{document}
